\documentclass[aspectratio=169,10pt,spanish]{beamer}

\input{../../../_preambulo_beamer.tex}
\input{../../../_comandos_beamer.tex}

\selectlanguage{spanish}

\title{MA1001B - Modelación Estadística para la Toma de Decisiones}
\subtitle{Lección 19: Variables aleatorias continuas: densidad y probabilidad de intervalo}
\author{}
\institute{Tecnol\'ogico de Monterrey}
\date{}

\begin{document}

\begin{frame}[plain]
  \vspace{1.2cm}
  {\Large\bfseries MA1001B - Modelación Estadística para la Toma de Decisiones\par}
  \vspace{0.7cm}
  {\large Lección 19: Variables aleatorias continuas\par}
  \vspace{0.7cm}
  {\color{TecRojo}\rule{\textwidth}{1.2pt}}
  \vspace{0.8cm}

  Facultad MA1001B\\[0.3cm]
  Periodo\\[0.3cm]
  Tecnol\'ogico de Monterrey
\end{frame}

\begin{frame}{La pregunta de probabilidad continua}
  \begin{alertblock}{Pregunta guía}
    Si el tiempo de servicio puede tomar cualquier valor en $[2, 14]$ minutos,
    ¿dónde vive la probabilidad: en un punto o en un intervalo?
  \end{alertblock}

  \vspace{0.6em}
  Al final de esta lección, usted puede:
  \begin{itemize}
    \item distinguir una densidad $f(x)$ de una FMP discreta;
    \item leer $P(a<X<b)$ como área bajo $f$;
    \item calcular esa área a partir de $F(b)-F(a)$;
    \item explicar por qué $P(X=c)=0$ para $X$ continua.
  \end{itemize}

  \vfill
  \scriptsize Todos los tiempos de servicio usados hoy son sintéticos. No se usan datos reales de empresa.
\end{frame}

\begin{frame}{Una densidad triangular de tiempo de servicio}
  \small
  Soporte $X\in[2,14]$ minutos, moda $6$. La altura pico es
  \[
    f(6)=\frac{2}{14-2}=\frac{1}{6}=0.166667.
  \]
  La media de una densidad triangular es
  \[
    E[X]=\frac{2+14+6}{3}=7.333333.
  \]

  \vspace{0.4em}
  \begin{exampleblock}{Densidad frente a probabilidad}
    $f(6)$ es una altura. La probabilidad es el área bajo $f$. El área total
    en $[2,14]$ es $1$.
  \end{exampleblock}
\end{frame}

\begin{frame}[fragile]{Paso 1: La densidad es una altura}
  \begin{columns}[T]
    \begin{column}{0.42\textwidth}
      \small
      \begin{lstlisting}
model = triang(
  c=1/3, loc=2, scale=12
)
f_mode = model.pdf(6)
area = model.cdf(14) - model.cdf(2)
      \end{lstlisting}

      \begin{block}{Salida verificada}
        $f(2)=0.000000$\\
        $f(6)=0.166667$\\
        Media $=7.333333$\\
        Área total $=1.000000$
      \end{block}
    \end{column}
    \begin{column}{0.58\textwidth}
      \centering
      \includegraphics[width=\textwidth]{../figuras/va_continua_01_densidad.png}
      \vspace{0.2em}
      \scriptsize Densidad triangular del Paso 1
    \end{column}
  \end{columns}
\end{frame}

\begin{frame}{La probabilidad de intervalo es área}
  \small
  Para una variable aleatoria continua,
  \[
    P(a<X<b)=F(b)-F(a)=\int_a^b f(x)\,dx.
  \]
  En el reloj sintético,
  \[
    F(4)=0.083333,\qquad F(10)=0.833333,
  \]
  \[
    P(4<X<10)=0.833333-0.083333=0.750000.
  \]
  El intervalo cerrado $P(4\le X\le 10)$ devuelve el mismo valor.
\end{frame}

\begin{frame}[fragile]{Paso 2: Sombrear $P(4<X<10)$}
  \begin{columns}[T]
    \begin{column}{0.40\textwidth}
      \small
      \begin{lstlisting}
area = (
    model.cdf(10)
    - model.cdf(4)
)
      \end{lstlisting}

      \begin{block}{Salida verificada}
        $F(4)=0.083333$\\
        $F(10)=0.833333$\\
        $P(4<X<10)=0.750000$\\
        $P(4\le X\le 10)=0.750000$
      \end{block}
    \end{column}
    \begin{column}{0.60\textwidth}
      \centering
      \includegraphics[width=\textwidth]{../figuras/va_continua_02_probabilidad_intervalo.png}
      \vspace{0.2em}
      \scriptsize Intervalo sombreado del Paso 2
    \end{column}
  \end{columns}
\end{frame}

\begin{frame}{Paso 3: Un punto carga probabilidad 0}
  \begin{columns}[T]
    \begin{column}{0.42\textwidth}
      \small
      $P(|X-6|<0.50)=0.158854$\\
      $P(|X-6|<0.10)=0.033021$\\
      $P(|X-6|<0.01)=0.003330$

      \vspace{0.4em}
      \[
        F(6)=0.333333,\quad P(X=6)=0.
      \]

      \vspace{0.4em}
      \textbf{Límite:} la FDA no tiene salto, de modo que un minuto
      individual tiene masa 0. La Lección 20 usa una densidad uniforme plana
      y Monte Carlo.
    \end{column}
    \begin{column}{0.58\textwidth}
      \centering
      \includegraphics[width=\textwidth]{../figuras/va_continua_03_limite_masa_puntual.png}
      \vspace{0.2em}
      \scriptsize FDA y ventanas que se encogen del Paso 3
    \end{column}
  \end{columns}
\end{frame}

\begin{frame}{Verificación de conceptos}
  \begin{enumerate}
    \item ¿Por qué $f(6)=0.166667$ no es igual a $P(X=6)$?
    \item Calcule $P(4<X<10)$ a partir de $F(10)=0.833333$ y $F(4)=0.083333$.
    \item ¿Por qué coinciden $P(4<X<10)$ y $P(4\le X\le 10)$?
    \item ¿Qué ocurre con $P(|X-6|<h)$ cuando $h\to 0$?
  \end{enumerate}

  \vspace{0.6em}
  \begin{block}{Conexión con la Lección 20}
    La sesión puede continuar directamente con el modelo uniforme continuo y
    una comprobación Monte Carlo de la probabilidad de intervalo.
  \end{block}

  \vfill
  \scriptsize Fuente: OpenStax, \textit{Introductory Business Statistics 2e},
  Capítulo 5, Sección 5.1. CC BY-NC-SA.
\end{frame}

\appendix



\end{document}
